\section{Analysis}
\label{sec:analysis}

\subsection{Structural Sources of Large Parallel Components}
\label{sec:why}

The prevalence of the parallel component is not only a learned artifact. It is
also encouraged by the pre-norm residual design. At layer $l$, the residual
norm can accumulate with depth:
\begin{equation}
  \|\hl{l}\| \approx \|\hl{0}\| + \sum_{k<l} \|\Dl{k}\|,
\end{equation}
while the block output $\Dl{l} = f^l(\LN(\hl{l}))$ is computed on a
\emph{normalized} input with $\|\LN(\hl{l})\| \approx \sqrt{d}$,
so $\|\Dl{l}\|$ need not grow with the residual norm. As $\|\hl{l}\|$ grows
and $\|\Dl{l}\|$ remains comparatively bounded, the residual update becomes a
smaller perturbation of a large vector. The $\LN$ operation removes the
magnitude of $\hl{l}$ before passing it to $f^l$, but it preserves directional
information. Learned block outputs can therefore remain correlated with the
current residual direction even when their norms are controlled. This provides
a structural reason to expect non-negligible parallel components across depth
and block type.

As a simple reference model, suppose $\Dl{l}$ is a random vector with fixed
magnitude $\epsilon$ and a direction drawn uniformly from the unit sphere in
$\mathbb{R}^d$. Under this model, the expected parallel ratio is:
\begin{equation}
  \mathbb{E}[\pr{l}{t}]
    = \mathbb{E}\!\left[\frac{|\Dl{l} \cdot \hat{\hl{l}}|}{\|\Dl{l}\|}\right]
    \approx \frac{1}{\sqrt{d}},
\end{equation}
\noindent
which gives a natural baseline before any learned correlation with the hidden
state is taken into account. Trained models can exceed this baseline when
$f^l$'s output direction remains correlated with $\hl{l}$ through the
normalization path. This toy-model argument is not intended as a strict theorem
for trained transformers, but as a structural explanation for why substantial
$\PR{l}$ is plausible under pre-norm residual accumulation.

This argument applies equally to $\Attn$ and $\FFN$ blocks, explaining the
universality of the phenomenon observed in Figure~\ref{fig:profiles}.

\subsection{Parallel Component and Attention Diagonal: Empirical Validation}
\label{sec:attn-diagonal-analysis}

Section~\ref{sec:value-level} predicts that one source of the attention-side
parallel component is the product of diagonal self-routing and value
self-alignment. Empirically, across layers and models, the attention-only
parallel ratio $\PR{l}$ tracks the mean diagonal attention weight
$\mathbb{E}_{x,t}[\Adiag{t}]$ closely.

This relationship supports the factorization in
Eq.~\ref{eq:alpha-attn}: layers where tokens route more attention to
themselves tend to produce larger attention-side parallel updates. The
diagonal term is not the only possible source of parallel residual geometry,
but it provides a measurable attention-side mechanism.
